\documentclass[12pt, a4paper]{article}

\usepackage[a4paper, margin=1in]{geometry}

\usepackage[utf8]{inputenc}
\usepackage[mathscr]{euscript}
\usepackage{amssymb,amsmath,amsthm,amsfonts}
\usepackage[shortlabels]{enumitem}
\usepackage{multicol,multirow}
\usepackage{lipsum}
\usepackage{balance}
\usepackage{calc}
\usepackage[colorlinks=true,citecolor=blue,linkcolor=blue]{hyperref}
\usepackage{import}
\usepackage{xifthen}
\usepackage{pdfpages}
\usepackage{transparent}
\usepackage{listings}

\newcommand{\incfig}[2][1.0]{
    \def\svgwidth{#1\columnwidth}
    \import{./figures/}{#2.pdf_tex}
}

\newcommand{\incimg}[2][1.0]{
  \includegraphics[width=#1\columnwidth]{./figures/#2}
}

\newlist{enumproof}{enumerate}{4}
\setlist[enumproof,1]{label=\arabic*., parsep=1em}
\setlist[enumproof,2]{label=\arabic{enumproofi}.\arabic*., parsep=1em}
\setlist[enumproof,3]{label=\arabic{enumproofi}.\arabic{enumproofii}.\arabic*., parsep=1em}
\setlist[enumproof,4]{label=\arabic{enumproofi}.\arabic{enumproofii}.\arabic{enumproofiii}.\arabic*., parsep=1em}

\renewcommand{\qedsymbol}{\ensuremath{\blacksquare}}

\lstdefinestyle{mystyle}{
  language=Python, % Set the language to Python
  commentstyle=\color{codegray}\itshape, % Python comments (#)
  keywordstyle=\color{orange}\bfseries,  % Python reserved keywords
  stringstyle=\color{mauve},             % Strings
  basicstyle={\ttfamily\footnotesize},   % Monospace font
  breakatwhitespace=true,         
  breaklines=true,                 
  captionpos=b,                    
  keepspaces=true,                 
  numbers=none,                    
  tabsize=2,
  showstringspaces=false,
  % Extra keyword groups for Python
  morekeywords=[2]{self, cls, __init__, __str__, __repr__}, 
  keywordstyle=[2]\color{identblue},     % Special identifiers / methods
  morekeywords=[3]{int, float, bool, str, list, dict, set, tuple, object, None, True, False},
  keywordstyle=[3]\color{codegreen},     % Built-in types
  morekeywords=[4]{print, input, open, range, len, type, isinstance, super, enumerate, zip, map, filter, sum, min, max, sorted},
  keywordstyle=[4]\color{cyan},          % Built-in functions
}

% Things Lie
\newcommand{\kb}{\mathfrak b}
\newcommand{\kg}{\mathfrak g}
\newcommand{\kh}{\mathfrak h}
\newcommand{\kn}{\mathfrak n}
\newcommand{\ku}{\mathfrak u}
\newcommand{\kz}{\mathfrak z}
\DeclareMathOperator{\Ext}{Ext} % Ext functor
\DeclareMathOperator{\Tor}{Tor} % Tor functor
\newcommand{\gl}{\opname{\mathfrak{gl}}} % frak gl group
\renewcommand{\sl}{\opname{\mathfrak{sl}}} % frak sl group chktex 6

% More script letters etc.
\newcommand{\SA}{\mathcal A}
\newcommand{\SB}{\mathcal B}
\newcommand{\SC}{\mathcal C}
\newcommand{\SF}{\mathcal F}
\newcommand{\SG}{\mathcal G}
\newcommand{\SH}{\mathcal H}
\newcommand{\OO}{\mathcal O}

\newcommand{\SCA}{\mathscr A}
\newcommand{\SCB}{\mathscr B}
\newcommand{\SCC}{\mathscr C}
\newcommand{\SCD}{\mathscr D}
\newcommand{\SCE}{\mathscr E}
\newcommand{\SCF}{\mathscr F}
\newcommand{\SCG}{\mathscr G}
\newcommand{\SCH}{\mathscr H}

% Mathfrak primes
\newcommand{\km}{\mathfrak m}
\newcommand{\kp}{\mathfrak p}
\newcommand{\kq}{\mathfrak q}

% number sets
\newcommand{\RR}[1][]{\ensuremath{\ifstrempty{#1}{\mathbb{R}}{\mathbb{R}^{#1}}}}
\newcommand{\NN}[1][]{\ensuremath{\ifstrempty{#1}{\mathbb{N}}{\mathbb{N}^{#1}}}}
\newcommand{\ZZ}[1][]{\ensuremath{\ifstrempty{#1}{\mathbb{Z}}{\mathbb{Z}^{#1}}}}
\newcommand{\QQ}[1][]{\ensuremath{\ifstrempty{#1}{\mathbb{Q}}{\mathbb{Q}^{#1}}}}
\newcommand{\CC}[1][]{\ensuremath{\ifstrempty{#1}{\mathbb{C}}{\mathbb{C}^{#1}}}}
\newcommand{\PP}[1][]{\ensuremath{\ifstrempty{#1}{\mathbb{P}}{\mathbb{P}^{#1}}}}
\newcommand{\HH}[1][]{\ensuremath{\ifstrempty{#1}{\mathbb{H}}{\mathbb{H}^{#1}}}}
\newcommand{\FF}[1][]{\ensuremath{\ifstrempty{#1}{\mathbb{F}}{\mathbb{F}^{#1}}}}
% expected value
\newcommand{\EE}{\ensuremath{\mathbb{E}}}
\newcommand{\charin}{\text{ char }}
\DeclareMathOperator{\sign}{sign}
\DeclareMathOperator{\Aut}{Aut}
\DeclareMathOperator{\Inn}{Inn}
\DeclareMathOperator{\Syl}{Syl}
\DeclareMathOperator{\Gal}{Gal}
\DeclareMathOperator{\GL}{GL} % General linear group
\DeclareMathOperator{\SL}{SL} % Special linear group

%---------------------------------------
% BlackBoard Math Fonts :-
%---------------------------------------

%Captital Letters
\newcommand{\bbA}{\mathbb{A}}	\newcommand{\bbB}{\mathbb{B}}
\newcommand{\bbC}{\mathbb{C}}	\newcommand{\bbD}{\mathbb{D}}
\newcommand{\bbE}{\mathbb{E}}	\newcommand{\bbF}{\mathbb{F}}
\newcommand{\bbG}{\mathbb{G}}	\newcommand{\bbH}{\mathbb{H}}
\newcommand{\bbI}{\mathbb{I}}	\newcommand{\bbJ}{\mathbb{J}}
\newcommand{\bbK}{\mathbb{K}}	\newcommand{\bbL}{\mathbb{L}}
\newcommand{\bbM}{\mathbb{M}}	\newcommand{\bbN}{\mathbb{N}}
\newcommand{\bbO}{\mathbb{O}}	\newcommand{\bbP}{\mathbb{P}}
\newcommand{\bbQ}{\mathbb{Q}}	\newcommand{\bbR}{\mathbb{R}}
\newcommand{\bbS}{\mathbb{S}}	\newcommand{\bbT}{\mathbb{T}}
\newcommand{\bbU}{\mathbb{U}}	\newcommand{\bbV}{\mathbb{V}}
\newcommand{\bbW}{\mathbb{W}}	\newcommand{\bbX}{\mathbb{X}}
\newcommand{\bbY}{\mathbb{Y}}	\newcommand{\bbZ}{\mathbb{Z}}

%---------------------------------------
% MathCal Fonts :-
%---------------------------------------

%Captital Letters
\newcommand{\mcA}{\mathcal{A}}	\newcommand{\mcB}{\mathcal{B}}
\newcommand{\mcC}{\mathcal{C}}	\newcommand{\mcD}{\mathcal{D}}
\newcommand{\mcE}{\mathcal{E}}	\newcommand{\mcF}{\mathcal{F}}
\newcommand{\mcG}{\mathcal{G}}	\newcommand{\mcH}{\mathcal{H}}
\newcommand{\mcI}{\mathcal{I}}	\newcommand{\mcJ}{\mathcal{J}}
\newcommand{\mcK}{\mathcal{K}}	\newcommand{\mcL}{\mathcal{L}}
\newcommand{\mcM}{\mathcal{M}}	\newcommand{\mcN}{\mathcal{N}}
\newcommand{\mcO}{\mathcal{O}}	\newcommand{\mcP}{\mathcal{P}}
\newcommand{\mcQ}{\mathcal{Q}}	\newcommand{\mcR}{\mathcal{R}}
\newcommand{\mcS}{\mathcal{S}}	\newcommand{\mcT}{\mathcal{T}}
\newcommand{\mcU}{\mathcal{U}}	\newcommand{\mcV}{\mathcal{V}}
\newcommand{\mcW}{\mathcal{W}}	\newcommand{\mcX}{\mathcal{X}}
\newcommand{\mcY}{\mathcal{Y}}	\newcommand{\mcZ}{\mathcal{Z}}


%---------------------------------------
% Bold Math Fonts :-
%---------------------------------------

%Captital Letters
\newcommand{\bmA}{\boldsymbol{A}}	\newcommand{\bmB}{\boldsymbol{B}}
\newcommand{\bmC}{\boldsymbol{C}}	\newcommand{\bmD}{\boldsymbol{D}}
\newcommand{\bmE}{\boldsymbol{E}}	\newcommand{\bmF}{\boldsymbol{F}}
\newcommand{\bmG}{\boldsymbol{G}}	\newcommand{\bmH}{\boldsymbol{H}}
\newcommand{\bmI}{\boldsymbol{I}}	\newcommand{\bmJ}{\boldsymbol{J}}
\newcommand{\bmK}{\boldsymbol{K}}	\newcommand{\bmL}{\boldsymbol{L}}
\newcommand{\bmM}{\boldsymbol{M}}	\newcommand{\bmN}{\boldsymbol{N}}
\newcommand{\bmO}{\boldsymbol{O}}	\newcommand{\bmP}{\boldsymbol{P}}
\newcommand{\bmQ}{\boldsymbol{Q}}	\newcommand{\bmR}{\boldsymbol{R}}
\newcommand{\bmS}{\boldsymbol{S}}	\newcommand{\bmT}{\boldsymbol{T}}
\newcommand{\bmU}{\boldsymbol{U}}	\newcommand{\bmV}{\boldsymbol{V}}
\newcommand{\bmW}{\boldsymbol{W}}	\newcommand{\bmX}{\boldsymbol{X}}
\newcommand{\bmY}{\boldsymbol{Y}}	\newcommand{\bmZ}{\boldsymbol{Z}}
%Small Letters
\newcommand{\bma}{\boldsymbol{a}}	\newcommand{\bmb}{\boldsymbol{b}}
\newcommand{\bmc}{\boldsymbol{c}}	\newcommand{\bmd}{\boldsymbol{d}}
\newcommand{\bme}{\boldsymbol{e}}	\newcommand{\bmf}{\boldsymbol{f}}
\newcommand{\bmg}{\boldsymbol{g}}	\newcommand{\bmh}{\boldsymbol{h}}
\newcommand{\bmi}{\boldsymbol{i}}	\newcommand{\bmj}{\boldsymbol{j}}
\newcommand{\bmk}{\boldsymbol{k}}	\newcommand{\bml}{\boldsymbol{l}}
\newcommand{\bmm}{\boldsymbol{m}}	\newcommand{\bmn}{\boldsymbol{n}}
\newcommand{\bmo}{\boldsymbol{o}}	\newcommand{\bmp}{\boldsymbol{p}}
\newcommand{\bmq}{\boldsymbol{q}}	\newcommand{\bmr}{\boldsymbol{r}}
\newcommand{\bms}{\boldsymbol{s}}	\newcommand{\bmt}{\boldsymbol{t}}
\newcommand{\bmu}{\boldsymbol{u}}	\newcommand{\bmv}{\boldsymbol{v}}
\newcommand{\bmw}{\boldsymbol{w}}	\newcommand{\bmx}{\boldsymbol{x}}
\newcommand{\bmy}{\boldsymbol{y}}	\newcommand{\bmz}{\boldsymbol{z}}

%---------------------------------------
% Scr Math Fonts :-
%---------------------------------------

\newcommand{\sA}{{\mathscr{A}}}   \newcommand{\sB}{{\mathscr{B}}}
\newcommand{\sC}{{\mathscr{C}}}   \newcommand{\sD}{{\mathscr{D}}}
\newcommand{\sE}{{\mathscr{E}}}   \newcommand{\sF}{{\mathscr{F}}}
\newcommand{\sG}{{\mathscr{G}}}   \newcommand{\sH}{{\mathscr{H}}}
\newcommand{\sI}{{\mathscr{I}}}   \newcommand{\sJ}{{\mathscr{J}}}
\newcommand{\sK}{{\mathscr{K}}}   \newcommand{\sL}{{\mathscr{L}}}
\newcommand{\sM}{{\mathscr{M}}}   \newcommand{\sN}{{\mathscr{N}}}
\newcommand{\sO}{{\mathscr{O}}}   \newcommand{\sP}{{\mathscr{P}}}
\newcommand{\sQ}{{\mathscr{Q}}}   \newcommand{\sR}{{\mathscr{R}}}
\newcommand{\sS}{{\mathscr{S}}}   \newcommand{\sT}{{\mathscr{T}}}
\newcommand{\sU}{{\mathscr{U}}}   \newcommand{\sV}{{\mathscr{V}}}
\newcommand{\sW}{{\mathscr{W}}}   \newcommand{\sX}{{\mathscr{X}}}
\newcommand{\sY}{{\mathscr{Y}}}   \newcommand{\sZ}{{\mathscr{Z}}}


%---------------------------------------
% Math Fraktur Font
%---------------------------------------

%Captital Letters
\newcommand{\mfA}{\mathfrak{A}}	\newcommand{\mfB}{\mathfrak{B}}
\newcommand{\mfC}{\mathfrak{C}}	\newcommand{\mfD}{\mathfrak{D}}
\newcommand{\mfE}{\mathfrak{E}}	\newcommand{\mfF}{\mathfrak{F}}
\newcommand{\mfG}{\mathfrak{G}}	\newcommand{\mfH}{\mathfrak{H}}
\newcommand{\mfI}{\mathfrak{I}}	\newcommand{\mfJ}{\mathfrak{J}}
\newcommand{\mfK}{\mathfrak{K}}	\newcommand{\mfL}{\mathfrak{L}}
\newcommand{\mfM}{\mathfrak{M}}	\newcommand{\mfN}{\mathfrak{N}}
\newcommand{\mfO}{\mathfrak{O}}	\newcommand{\mfP}{\mathfrak{P}}
\newcommand{\mfQ}{\mathfrak{Q}}	\newcommand{\mfR}{\mathfrak{R}}
\newcommand{\mfS}{\mathfrak{S}}	\newcommand{\mfT}{\mathfrak{T}}
\newcommand{\mfU}{\mathfrak{U}}	\newcommand{\mfV}{\mathfrak{V}}
\newcommand{\mfW}{\mathfrak{W}}	\newcommand{\mfX}{\mathfrak{X}}
\newcommand{\mfY}{\mathfrak{Y}}	\newcommand{\mfZ}{\mathfrak{Z}}
%Small Letters
\newcommand{\mfa}{\mathfrak{a}}	\newcommand{\mfb}{\mathfrak{b}}
\newcommand{\mfc}{\mathfrak{c}}	\newcommand{\mfd}{\mathfrak{d}}
\newcommand{\mfe}{\mathfrak{e}}	\newcommand{\mff}{\mathfrak{f}}
\newcommand{\mfg}{\mathfrak{g}}	\newcommand{\mfh}{\mathfrak{h}}
\newcommand{\mfi}{\mathfrak{i}}	\newcommand{\mfj}{\mathfrak{j}}
\newcommand{\mfk}{\mathfrak{k}}	\newcommand{\mfl}{\mathfrak{l}}
\newcommand{\mfm}{\mathfrak{m}}	\newcommand{\mfn}{\mathfrak{n}}
\newcommand{\mfo}{\mathfrak{o}}	\newcommand{\mfp}{\mathfrak{p}}
\newcommand{\mfq}{\mathfrak{q}}	\newcommand{\mfr}{\mathfrak{r}}
\newcommand{\mfs}{\mathfrak{s}}	\newcommand{\mft}{\mathfrak{t}}
\newcommand{\mfu}{\mathfrak{u}}	\newcommand{\mfv}{\mathfrak{v}}
\newcommand{\mfw}{\mathfrak{w}}	\newcommand{\mfx}{\mathfrak{x}}
\newcommand{\mfy}{\mathfrak{y}}	\newcommand{\mfz}{\mathfrak{z}}


\newcommand{\mytitle}{FMA\_3TC13\_TP Poly Exercices 4}
\newcommand{\myauthor}{github/omgeta}
\newcommand{\mydate}{AU 25-26 S1}

\begin{document}
\raggedright
\footnotesize
\begin{center}
{\normalsize{\textbf{\mytitle}}} \\
{\footnotesize{\mydate\hspace{2pt}\textemdash\hspace{2pt}\myauthor}}
\end{center}
\setlist{topsep=-1em, itemsep=-1em, parsep=2em}
\setlist[enumerate]{topsep=0.4em,itemsep=0.8em,parsep=0.35em}
%%%%%%%%%%%%%%%%%%%%%%%%%%%%%%%%%%%%%%%%%%%%%%%%%%%%%%
%                      Begin                         %
%%%%%%%%%%%%%%%%%%%%%%%%%%%%%%%%%%%%%%%%%%%%%%%%%%%%%%
\begin{enumerate}[label=4.\arabic*., start=60]

  \item \begin{minipage}[t]{\linewidth}\raggedright
\textbf{Question.} Calculer l’inverse de $\displaystyle P=\begin{pmatrix}1&2&-1&3&1&2&4\\0&1&2&-1&3&1&2\\0&0&1&2&-1&3&1\\0&0&0&1&2&-1&3\\0&0&0&0&1&2&-1\\0&0&0&0&0&1&2\\0&0&0&0&0&0&1\end{pmatrix}$.
\par\smallskip
\textbf{Réponse.} La réduction de $[P\mid I_7]$ donne $[I_7\mid P^{-1}]$, avec
    \[
      P^{-1}=\begin{pmatrix}
        1&-2&5&-15&40&-110&300\\
        0&1&-2&5&-15&40&-110\\
        0&0&1&-2&5&-15&40\\
        0&0&0&1&-2&5&-15\\
        0&0&0&0&1&-2&5\\
        0&0&0&0&0&1&-2\\
        0&0&0&0&0&0&1
      \end{pmatrix}
    \]
    Le produit $PP^{-1}=I_7$ vérifie le résultat.
\end{minipage}

  \item \begin{minipage}[t]{\linewidth}\raggedright
\textbf{Question.} Existe-t-il $X\in\RR^{3\times3}$ telle que $XA=B$, avec $\displaystyle A=\begin{pmatrix}1&2&-1\\2&-1&-1\\-5&0&3\end{pmatrix}$ et $\displaystyle B=\begin{pmatrix}-2&1&1\\8&-1&-5\\4&3&-3\end{pmatrix}$ ?
\par\smallskip
\textbf{Réponse.} Non. Pour $v=(3,1,5)^\top$, on calcule $Av=0$ mais $Bv=(0,-2,0)^\top\neq0$.
    Si $XA=B$, alors $Bv=XAv=0$, contradiction.
\end{minipage}

  \item \begin{minipage}[t]{\linewidth}\raggedright
\textbf{Question.} Soient $n\geq1$, $\zeta=e^{2i\pi/n}$, $A_{i,j}=\zeta^{(i-1)(j-1)}$ et $B_{i,j}=\zeta^{(i-1)(1-j)}$ pour $1\leq i,j\leq n$. Calculer $A^2,B^2,AB,BA,A^{-1},B^{-1}$.
\par\smallskip
\textbf{Réponse.} Posons $R_{i,j}=1$ si $i+j\equiv2\pmod n$, et $R_{i,j}=0$ sinon.
    La somme géométrique donne
    \[
      \sum_{k=0}^{n-1}\zeta^{k\ell}=\begin{cases}n&n\mid\ell,\\0&n\nmid\ell.\end{cases}
    \]
    Ainsi $(A^2)_{i,j}=\sum_{k=0}^{n-1}\zeta^{k(i+j-2)}$ et
    $(AB)_{i,j}=\sum_{k=0}^{n-1}\zeta^{k(i-j)}$.
    De même pour $B^2$ et $BA$,
    \[
      \therefore A^2=B^2=nR,\quad AB=BA=nI_n,\quad A^{-1}=\frac1nB,\quad B^{-1}=\frac1nA
    \]
\end{minipage}

  \item \begin{minipage}[t]{\linewidth}\raggedright
\textbf{Énoncé.} On pose $u_{n+1}=(3u_n+2)/(u_n+2)$, $v_0=u_0$, $w_0=1$, $v_{n+1}=3v_n+2w_n$ et $w_{n+1}=v_n+2w_n$.

\begin{enumerate}[\arabic*.]
        \item \textbf{Question.} Montrer que $u_n=v_n/w_n$.
\par\smallskip
\textbf{Réponse.} Tant que la suite est définie, $u_n=v_n/w_n$ par récurrence : c'est vrai pour $n=0$, et
        \[
          \frac{v_{n+1}}{w_{n+1}}=\frac{3v_n+2w_n}{v_n+2w_n}
          =\frac{3u_n+2}{u_n+2}=u_{n+1}
        \]

  \item \textbf{Question.} Trouver $A\in\RR^{2\times2}$ telle que $(v_{n+1},w_{n+1})^\top=A(v_n,w_n)^\top$.
\par\smallskip
\textbf{Réponse.} $\displaystyle A=\begin{pmatrix}3&2\\1&2\end{pmatrix}$.

  \item \textbf{Question.} Pour $P=\begin{pmatrix}-1&2\\1&1\end{pmatrix}$, calculer $P^{-1}AP$ et en déduire $A^n$.
\par\smallskip
\textbf{Réponse.} $\displaystyle P^{-1}=\frac13\begin{pmatrix}-1&2\\1&1\end{pmatrix}$ et
        $\displaystyle P^{-1}AP=\begin{pmatrix}1&0\\0&4\end{pmatrix}$, donc
        \[
          A^n=\frac13\begin{pmatrix}1+2\cdot4^n&2\cdot4^n-2\\4^n-1&4^n+2\end{pmatrix}
        \]

  \item \textbf{Question.} Calculer $v_n,w_n$ et $\lim_{n\to\infty}u_n$.
\par\smallskip
\textbf{Réponse.} En appliquant $A^n$ à $(u_0,1)^\top$,
        \[
          v_n=\frac{2(u_0+1)4^n+u_0-2}{3},\qquad
          w_n=\frac{(u_0+1)4^n-u_0+2}{3}
        \]
        Si $u_0=-1$, $u_n=-1$ pour tout $n$ et la limite vaut $-1$.
        Sinon, dès que la suite est définie pour tout $n$, $\lim u_n=2$.
        Les valeurs initiales exclues sont $u_0=-(4^r+2)/(4^r-1)$ pour $r\geq1$ : elles donnent $w_r=0$.

\end{enumerate}
\end{minipage}

  \item \begin{minipage}[t]{\linewidth}\raggedright
\textbf{Question.} Soient $a,b\in K$ et $M\in K^{n\times n}$, de diagonale $a$, première surdiagonale $b$, et tous les autres coefficients nuls. Calculer $M^k$ pour tout $k\in\NN$.
\par\smallskip
\textbf{Réponse.} Soit $N$ la matrice dont la première surdiagonale vaut $1$ et les autres coefficients $0$.
    Alors $M=aI_n+bN$, $N^n=0$, et
    \[
      \therefore M^k=\sum_{r=0}^{\min(k,n-1)}\binom kr a^{k-r}b^rN^r
    \]
    Autrement dit, $(M^k)_{i,j}=\binom{k}{j-i}a^{k-j+i}b^{j-i}$ si $0\leq j-i\leq k$, et $0$ sinon.
    On prend $M^0=I_n$.
\end{minipage}

  \item \begin{minipage}[t]{\linewidth}\raggedright
\textbf{Question.} Calculer les puissances de $K\in\QQ^{n\times n}$ définie par $K_{1,j}=K_{j,1}=1$ pour $j>1$, tous les autres coefficients étant nuls.
\par\smallskip
\textbf{Réponse.} En notant $J_{n-1}$ la matrice pleine de $1$,
    \[
      K^2=\begin{pmatrix}n-1&0\\0&J_{n-1}\end{pmatrix},\qquad K^3=(n-1)K
    \]
    Donc $K^0=I_n$ et, pour $r\geq0$,
    \[
      K^{2r+1}=(n-1)^rK,\qquad K^{2r+2}=(n-1)^rK^2
    \]
\end{minipage}

  \item \begin{minipage}[t]{\linewidth}\raggedright
\textbf{Énoncé.} Pour $\alpha\in\RR$, discuter l’inversibilité et calculer l’inverse lorsqu’il existe.

\begin{enumerate}[\arabic*.]
        \item \textbf{Question.} $\displaystyle A=\begin{pmatrix}\cosh\alpha&\sinh\alpha\\\sinh\alpha&\cosh\alpha\end{pmatrix}$.
\par\smallskip
\textbf{Réponse.} $\det A=\cosh^2\alpha-\sinh^2\alpha=1$.
        $\displaystyle\therefore A^{-1}=\begin{pmatrix}\cosh\alpha&-\sinh\alpha\\-\sinh\alpha&\cosh\alpha\end{pmatrix}$.

  \item \textbf{Question.} $\displaystyle B=\begin{pmatrix}\alpha&1&1\\1&\alpha&1\\1&1&\alpha\end{pmatrix}$.
\par\smallskip
\textbf{Réponse.} $B=(\alpha-1)I_3+J_3$, où $J_3^2=3J_3$.
        Alors $\det B=(\alpha-1)^2(\alpha+2)$.
        Si $\alpha\notin\{1,-2\}$,
        \[
          B^{-1}=\frac1{\alpha-1}I_3-\frac1{(\alpha-1)(\alpha+2)}J_3
        \]
        Sinon $B$ n'est pas inversible.

  \item \textbf{Question.} $\displaystyle C=\begin{pmatrix}1&2&1\\1&1&\alpha+1\\1&1&-\alpha^2\end{pmatrix}$.
\par\smallskip
\textbf{Réponse.} Posons $d=\alpha^2+\alpha+1=(\alpha+1/2)^2+3/4>0$.
        $\det C=d$, donc $C$ est toujours inversible sur $\RR$, avec
        \[
          C^{-1}=\frac1d\begin{pmatrix}-d&2\alpha^2+1&2\alpha+1\\d&-\alpha^2-1&-\alpha\\0&1&-1\end{pmatrix}
        \]

\end{enumerate}
\end{minipage}

  \item \begin{minipage}[t]{\linewidth}\raggedright
\textbf{Énoncé.} On considère $\displaystyle A=\begin{pmatrix}0&-1&1\\1&2&-3\\1&1&-2\end{pmatrix}$ et $\displaystyle P=\begin{pmatrix}1&1&0\\1&-1&1\\1&0&1\end{pmatrix}$.

\begin{enumerate}[\arabic*.]
        \item \textbf{Question.} Montrer que $P$ est inversible et calculer $P^{-1}$.
\par\smallskip
\textbf{Réponse.} $\det P=-1\neq0$ et $\displaystyle P^{-1}=\begin{pmatrix}1&1&-1\\0&-1&1\\-1&-1&2\end{pmatrix}$.

  \item \textbf{Question.} Calculer $D=P^{-1}AP$.
\par\smallskip
\textbf{Réponse.} $D=P^{-1}AP=\operatorname{diag}(0,1,-1)$.

  \item \textbf{Question.} Calculer $D^n$ pour tout $n\in\NN$.
\par\smallskip
\textbf{Réponse.} $D^0=I_3$ et $D^n=\operatorname{diag}(0,1,(-1)^n)$ pour $n\geq1$.

  \item \textbf{Question.} En déduire $A^n$ pour tout $n\in\NN$.
\par\smallskip
\textbf{Réponse.} $A^0=I_3$ et, pour $n\geq1$, en posant $s=(-1)^n$,
        \[
          A^n=PD^nP^{-1}=\begin{pmatrix}0&-1&1\\-s&-1-s&1+2s\\-s&-s&2s\end{pmatrix}
        \]

\end{enumerate}
\end{minipage}

  \item \begin{minipage}[t]{\linewidth}\raggedright
\textbf{Question.} Soient $x,y,z\in K$. Calculer $A^k$ pour tout $k\in\NN$, pour la matrice imprimée $\displaystyle A=\begin{pmatrix}x^2&xy&xz\\xy&y^2&yz\\xz&xy&y^2\end{pmatrix}$.
\par\smallskip
\textbf{Réponse.} Le corrigé utilise la matrice symétrique $A=vv^\top$, $v=(x,y,z)^\top$ : la dernière ligne est alors $(xz,yz,z^2)$.
    Pour cette version, $A^2=(v^\top v)A$, donc
    \[
      A^0=I_3,\qquad A^k=(x^2+y^2+z^2)^{k-1}A\quad(k\geq1)
    \]
    Pour la matrice imprimée, dont la dernière ligne est $(xz,xy,y^2)$, on a exactement
    \[
      A=UV,\quad U=\begin{pmatrix}x&0\\y&0\\0&1\end{pmatrix},\quad
      V=\begin{pmatrix}x&y&z\\xz&xy&y^2\end{pmatrix},\quad
      C=VU=\begin{pmatrix}x^2+y^2&z\\x^2z+xy^2&y^2\end{pmatrix}
    \]
    Ainsi $A^k=UC^{k-1}V$ pour $k\geq1$.
    Les puissances de $C$ se calculent par $C^0=I_2$, $C^1=C$ et
    $C^{r+2}=(x^2+2y^2)C^{r+1}-\det(C)C^r$.
\end{minipage}

  \item \begin{minipage}[t]{\linewidth}\raggedright
\textbf{Énoncé.} Soient $A,B\in\RR^{n\times n}$.

\begin{enumerate}[\arabic*.]
        \item \textbf{Question.} Si $\operatorname{tr}(AA^\top)=0$, que peut-on dire de $A$ ?
\par\smallskip
\textbf{Réponse.} $\operatorname{tr}(AA^\top)=\sum_{i,j}a_{i,j}^2=0$.
        Une somme de carrés réels est nulle seulement si chaque terme est nul. $\therefore A=0$.

  \item \textbf{Question.} Si $\operatorname{tr}(XA)=\operatorname{tr}(XB)$ pour tout $X\in\RR^{n\times n}$, que peut-on dire de $A,B$ ?
\par\smallskip
\textbf{Réponse.} Pour $X=E_{j,i}$, $\operatorname{tr}(XA)=a_{i,j}$ et $\operatorname{tr}(XB)=b_{i,j}$.
        $\therefore A=B$.

\end{enumerate}
\end{minipage}

  \item \begin{minipage}[t]{\linewidth}\raggedright
\textbf{Énoncé.} Soient $A,B\in K^{n\times n}$ inversibles. Simplifier au maximum.

\begin{enumerate}[\arabic*.]
        \item \textbf{Question.} $(2I_n+B^{-1})B+A(B-A^{-1}+2A^{-1}B)-AB$.
\par\smallskip
\textbf{Réponse.} $(2I_n+B^{-1})B+A(B-A^{-1}+2A^{-1}B)-AB
        =2B+I_n+AB-I_n+2B-AB=4B$.

  \item \textbf{Question.} $A(7I_n+A^{-1})-B(5B^{-1}A+B^{-1}+A)+A(B-2I_n)$.
\par\smallskip
\textbf{Réponse.} $A(7I_n+A^{-1})-B(5B^{-1}A+B^{-1}+A)+A(B-2I_n)
        =7A+I_n-5A-I_n-BA+AB-2A=AB-BA$.

\end{enumerate}
\end{minipage}

  \item \begin{minipage}[t]{\linewidth}\raggedright
\textbf{Question.} Soient $h_1,\ldots,h_n\in K\setminus\{0\}$. La matrice $M$ a pour diagonale $1+1/h_i$ et pour autres coefficients $1$. Calculer son inverse lorsqu’il existe.
\par\smallskip
\textbf{Réponse.} Posons $H=\operatorname{diag}(h_1,\ldots,h_n)$, $h=(h_1,\ldots,h_n)^\top$ et $s=1+\sum_i h_i$.
    Alors $M=H^{-1}+\boldsymbol{1}\boldsymbol{1}^\top$.
    Si $s\neq0$, un produit direct donne
    \[
      \therefore M^{-1}=H-\frac{hh^\top}{s},\qquad
      (M^{-1})_{i,j}=h_i\delta_{i,j}-\frac{h_ih_j}{s}
    \]
    Si $s=0$, $Mh=s\boldsymbol{1}=0$ alors que $h\neq0$ : $M$ n'est pas inversible.
\end{minipage}

  \item \begin{minipage}[t]{\linewidth}\raggedright
\textbf{Question.} Soient $\alpha_1,\ldots,\alpha_n\neq0$. La matrice $M$ a pour seuls coefficients non nuls $M_{i,n+1-i}=\alpha_i$. Calculer son inverse lorsqu’il existe.
\par\smallskip
\textbf{Réponse.} Soit $R$ la matrice qui renverse l'ordre des coordonnées. Alors $R^2=I_n$ et
    $M=\operatorname{diag}(\alpha_1,\ldots,\alpha_n)R$.
    \[
      \therefore M^{-1}=R\operatorname{diag}(\alpha_1^{-1},\ldots,\alpha_n^{-1})
    \]
    Ses seuls coefficients non nuls sont $(M^{-1})_{i,n+1-i}=1/\alpha_{n+1-i}$.
\end{minipage}

  \item \begin{minipage}[t]{\linewidth}\raggedright
\textbf{Énoncé.} Soient $\alpha,\beta,\gamma\in\RR$. On considère $x_1+2x_2-3x_3=\alpha$, $2x_1+6x_2-5x_3=\beta$, $x_1-2x_2+7x_3=\gamma$.

\begin{enumerate}[\arabic*.]
        \item \textbf{Question.} Écrire le système sous la forme $Ax=b$.
\par\smallskip
\textbf{Réponse.} $\displaystyle A=\begin{pmatrix}1&2&-3\\2&6&-5\\1&-2&7\end{pmatrix}$, $b=(\alpha,\beta,\gamma)^\top$.

  \item \textbf{Question.} Écrire et réduire la matrice augmentée ; en déduire les solutions.
\par\smallskip
\textbf{Réponse.} La réduction donne $[A\mid b]\sim[I_3\mid x]$, où
        \[
          x=\begin{pmatrix}(4\alpha-\beta+\gamma)/3\\(-19\alpha+10\beta-\gamma)/24\\(-5\alpha+2\beta+\gamma)/12\end{pmatrix}
        \]
        C'est l'unique solution.

  \item \textbf{Question.} En déduire $A^{-1}$.
\par\smallskip
\textbf{Réponse.} $\displaystyle A^{-1}=\frac1{24}\begin{pmatrix}32&-8&8\\-19&10&-1\\-10&4&2\end{pmatrix}$.

  \item \textbf{Question.} Décomposer $A=LU$ et résoudre à nouveau le système avec cette décomposition.
\par\smallskip
\textbf{Réponse.} $A=LU$, avec
        \[
          L=\begin{pmatrix}1&0&0\\2&1&0\\1&-2&1\end{pmatrix},\quad
          U=\begin{pmatrix}1&2&-3\\0&2&1\\0&0&12\end{pmatrix}
        \]
        On résout $Ly=b$ : $y=(\alpha,\beta-2\alpha,\gamma+2\beta-5\alpha)^\top$.
        Puis $Ux=y$ donne $x_3=y_3/12$, $x_2=(y_2-x_3)/2$, $x_1=y_1-2x_2+3x_3$, soit le résultat précédent.

\end{enumerate}
\end{minipage}

  \item \begin{minipage}[t]{\linewidth}\raggedright
\textbf{Énoncé.} Soient $a,b,c,d\in\RR$ et $\displaystyle A=\begin{pmatrix}a&a&a&a\\a&b&b&b\\a&b&c&c\\a&b&c&d\end{pmatrix}$.

\begin{enumerate}[\arabic*.]
        \item \textbf{Question.} Obtenir une décomposition $LU$ par élimination de Gauss.
\par\smallskip
\textbf{Réponse.} En soustrayant chaque ligne précédente, de bas en haut, on obtient
        \[
          A=LU,\quad L=\begin{pmatrix}1&0&0&0\\1&1&0&0\\1&1&1&0\\1&1&1&1\end{pmatrix},\quad
          U=\begin{pmatrix}a&a&a&a\\0&b-a&b-a&b-a\\0&0&c-b&c-b\\0&0&0&d-c\end{pmatrix}
        \]

  \item \textbf{Question.} Donner les conditions d’inversibilité de $A$.
\par\smallskip
\textbf{Réponse.} $\det A=a(b-a)(c-b)(d-c)$.
        $\therefore A$ est inversible si et seulement si $a\neq0$, $b\neq a$, $c\neq b$, $d\neq c$.

  \item \textbf{Question.} Résoudre $x_1+x_2+x_3+x_4=2$, $x_1+2x_2+2x_3+2x_4=1$, $x_1+2x_2+3x_3+3x_4=2$, $x_1+2x_2+3x_3+4x_4=3$.
\par\smallskip
\textbf{Réponse.} Les différences successives donnent
        $x_4=1$, $x_3+x_4=1$, $x_2+x_3+x_4=-1$, $x_1+x_2+x_3+x_4=2$.
        $\therefore (x_1,x_2,x_3,x_4)=(3,-2,0,1)$.

\end{enumerate}
\end{minipage}

  \item \begin{minipage}[t]{\linewidth}\raggedright
\textbf{Question.} Soient $\alpha\in K$, $A,B\in K^{n\times n}$. Résoudre en $X\in K^{n\times n}$ l’équation $\alpha X+\operatorname{tr}(X)A=B$.
\par\smallskip
\textbf{Réponse.} Posons $t=\operatorname{tr}X$. En prenant la trace,
    $(\alpha+\operatorname{tr}A)t=\operatorname{tr}B$.
    \begin{enumproof}
      \item Si $\alpha\neq0$ et $\alpha+\operatorname{tr}A\neq0$, l'unique solution est
        \[
          X=\frac1\alpha\left(B-\frac{\operatorname{tr}B}{\alpha+\operatorname{tr}A}A\right)
        \]
      \item Si $\alpha\neq0$ et $\alpha+\operatorname{tr}A=0$, il n'y a aucune solution si $\operatorname{tr}B\neq0$.
        Sinon, les solutions sont $X=(B-tA)/\alpha$, $t\in K$.
      \item Si $\alpha=0$ et $A\neq0$, il faut $B=tA$ pour un scalaire $t$.
        Si ce scalaire existe, les solutions sont toutes les matrices de trace $t$ ; sinon, aucune solution.
      \item Si $\alpha=0$ et $A=0$, toutes les matrices conviennent si $B=0$, et aucune sinon.
    \end{enumproof}
\end{minipage}

  \item \begin{minipage}[t]{\linewidth}\raggedright
\textbf{Question.} Soit $M=(m_{i,j})\in\RR^{n\times n}$ telle que $m_{i,i}>\sum_{j\neq i}m_{i,j}$ pour tout $i$. Montrer que $M$ est inversible (hypothèse telle qu’imprimée dans le poly).
\par\smallskip
\textbf{Réponse.} L'énoncé imprimé est faux sans valeurs absolues. Par exemple,
    $\displaystyle M=\begin{pmatrix}1&-1\\-1&1\end{pmatrix}$ vérifie $m_{i,i}>\sum_{j\neq i}m_{i,j}$, mais $\det M=0$.

    Avec l'hypothèse corrigée $|m_{i,i}|>\sum_{j\neq i}|m_{i,j}|$, supposons $Mx=0$ et $x\neq0$.
    Choisissons $i$ tel que $|x_i|=\max_j|x_j|>0$. Alors
    \[
      |m_{i,i}|\,|x_i|=\left|\sum_{j\neq i}m_{i,j}x_j\right|
      \leq\sum_{j\neq i}|m_{i,j}|\,|x_i|<|m_{i,i}|\,|x_i|
    \]
    Contradiction. $\therefore\ker M=\{0\}$, donc $M$ est inversible.
\end{minipage}

  \item \begin{minipage}[t]{\linewidth}\raggedright
\textbf{Question.} Soient $A\in\RR^{m\times n}$, $B\in\RR^{p\times q}$. Si $AXB=0$ pour tout $X\in\RR^{n\times p}$, montrer que $A=0$ ou $B=0$.
\par\smallskip
\textbf{Réponse.} Si $A\neq0$ et $B\neq0$, choisissons $a_{i,j}\neq0$ et $b_{k,\ell}\neq0$.
    Pour $X=E_{j,k}\in\RR^{n\times p}$, $(AXB)_{i,\ell}=a_{i,j}b_{k,\ell}\neq0$, contradiction.
    $\therefore A=0$ ou $B=0$.
\end{minipage}

  \item \begin{minipage}[t]{\linewidth}\raggedright
\textbf{Question.} Soit $A\in K^{n\times n}$ inversible telle que $A+A^{-1}=3I_n$. Calculer $A^3+A^{-3}$.
\par\smallskip
\textbf{Réponse.} $A$ et $A^{-1}$ commutent. Alors
    \[
      (A+A^{-1})^3=A^3+A^{-3}+3(A+A^{-1})=27I_n
    \]
    $\therefore A^3+A^{-3}=18I_n$.
\end{minipage}

  \item \begin{minipage}[t]{\linewidth}\raggedright
\textbf{Question.} Si $AB=BA$ et $A$ est inversible, montrer que $A^{-1}B=BA^{-1}$.
\par\smallskip
\textbf{Réponse.} $AB=BA\Rightarrow A^{-1}ABA^{-1}=A^{-1}BAA^{-1}\Rightarrow BA^{-1}=A^{-1}B$.
\end{minipage}

  \item \begin{minipage}[t]{\linewidth}\raggedright
\textbf{Question.} Pour $n\geq2$, déterminer les matrices de $K^{n\times n}$ qui commutent avec toutes les matrices symétriques.
\par\smallskip
\textbf{Réponse.} Si $M$ commute avec chaque $E_{i,i}$, ses coefficients hors diagonale sont nuls.
    Écrivons donc $M=\operatorname{diag}(d_1,\ldots,d_n)$.
    En commutant avec $E_{i,j}+E_{j,i}$, on obtient $d_i=d_j$.

    $\therefore M=\lambda I_n$, $\lambda\in K$.
    Réciproquement, toute matrice scalaire commute avec toutes les matrices.
\end{minipage}

  \item \begin{minipage}[t]{\linewidth}\raggedright
\textbf{Question.} Soit $A\in K^{n\times n}$ triangulaire supérieure, de diagonale $1$ et de coefficients $-1$ strictement au-dessus de la diagonale. Justifier son inversibilité et calculer son inverse.
\par\smallskip
\textbf{Réponse.} Écrivons $A=I_n-N$, où $N$ vaut $1$ strictement au-dessus de la diagonale et $0$ ailleurs.
    Comme $N^n=0$, $A^{-1}=I_n+N+\cdots+N^{n-1}$.
    Pour $i<j$, $(N^r)_{i,j}=\binom{j-i-1}{r-1}$ : on choisit les indices intermédiaires.
    \[
      \therefore (A^{-1})_{i,j}=\begin{cases}0&i>j,\\1&i=j,\\2^{j-i-1}&i<j.\end{cases}
    \]
\end{minipage}

  \item \begin{minipage}[t]{\linewidth}\raggedright
\textbf{Question.} Soit $\omega\in\RR$. La matrice triangulaire supérieure $M$ vérifie $M_{i,j}=\omega^{j-i}$ pour $j\geq i$. Justifier son inversibilité et calculer son inverse.
\par\smallskip
\textbf{Réponse.} Pour $N$ ayant des $1$ sur la première surdiagonale, $N^n=0$ et
    $M=I_n+\omega N+\cdots+\omega^{n-1}N^{n-1}$.
    Le produit télescopique donne $(I_n-\omega N)M=I_n$.

    $\therefore M^{-1}=I_n-\omega N$.
\end{minipage}

  \item \begin{minipage}[t]{\linewidth}\raggedright
\textbf{Énoncé.} Soient $U\in K^{m\times n}$ et $V\in K^{n\times m}$.

\begin{enumerate}[\arabic*.]
        \item \textbf{Question.} Supposer $I_m-UV$ inversible. Calculer $(I_n-VU)(I_n+V(I_m-UV)^{-1}U)$.
\par\smallskip
\textbf{Réponse.} Posons $C=(I_m-UV)^{-1}$. Comme $(I_n-VU)V=V(I_m-UV)$,
        \[
          AB=(I_n-VU)(I_n+VCU)=I_n-VU+V(I_m-UV)CU=I_n
        \]

  \item \textbf{Question.} Montrer que $I_m-UV$ est inversible si et seulement si $I_n-VU$ l’est.
\par\smallskip
\textbf{Réponse.} Les matrices étant carrées, $AB=I_n$ implique que $A$ est inversible.
        En échangeant $U$ et $V$, on obtient la réciproque.
        \[
          \therefore (I_n-VU)^{-1}=I_n+V(I_m-UV)^{-1}U
        \]

\end{enumerate}
\end{minipage}

  \item \begin{minipage}[t]{\linewidth}\raggedright
\textbf{Question.} Pour $M\in\RR^{n\times n}$, on pose $\sigma(M)=\sum_{i,j}m_{i,j}$. Étant donnée $A\in\RR^{n\times n}$, résoudre $X+X^\top=\sigma(X)A$.
\par\smallskip
\textbf{Réponse.} Posons $t=\sigma(X)$. L'équation donne $2t=t\sigma(A)$.
    Toute matrice antisymétrique $K$ est solution car $\sigma(K)=0$.

    Si $t\neq0$, $A=(X+X^\top)/t$ est symétrique et $\sigma(A)=2$.
    Dans ce cas, $X=(t/2)A+K$, avec $K^\top=-K$.
    Réciproquement, ces matrices vérifient l'équation.
    \[
      \therefore\mathcal S=\begin{cases}
        \{\lambda A+K:\lambda\in\RR,\ K^\top=-K\}&A=A^\top\text{ et }\sigma(A)=2,\\
        \{K:K^\top=-K\}&\text{sinon.}
      \end{cases}
    \]
\end{minipage}

\end{enumerate}
%%%%%%%%%%%%%%%%%%%%%%%%%%%%%%%%%%%%%%%%%%%%%%%%%%%%%%
%                       End                          %
%%%%%%%%%%%%%%%%%%%%%%%%%%%%%%%%%%%%%%%%%%%%%%%%%%%%%%

\end{document}
